\section{Literature Review}
\label{sec:lit}

\subsection{From Mean--Variance Optimisation to Risk-Based Allocation}

Modern portfolio construction begins with \citet{markowitz52}, who formalised
allocation as a trade-off between expected return and variance. The framework is
theoretically complete and practically fragile: its solutions are highly sensitive to
the expected-return vector, and sample covariance matrices estimated from short windows
on many assets are poorly conditioned in exactly the directions the optimiser exploits.
\citet{demiguel09} quantified the consequence, finding that a naive equal-weighted
portfolio outperformed a range of optimised alternatives out of sample once estimation
error was taken into account.

Two responses dominate. The first improves the inputs. \citet{ledoit04} introduced a
shrinkage estimator that trades a small bias for a large variance reduction and yields
a well-conditioned covariance matrix; the approach remains standard, and recent work
extends shrinkage to higher moments \citep{lu25}. The second abandons expected returns
altogether in favour of risk-based construction. Inverse-volatility weighting, risk
parity and hierarchical risk parity \citep{prado16} allocate on the basis of the risk
model alone, and \citet{asimit25} develop the risk-budgeting formulation further by
linking portfolio selection to risk sharing; \citet{khodamoradi23} pursue a related
robustness objective through multi-interval mean-conditional-value-at-risk
optimisation, and \citet{kehinde25} evaluate volatility targeting through an inverse
data-envelopment formulation. Volatility management is a related idea
applied at the portfolio level rather than the asset level: \citet{moreira17} showed
that scaling exposure inversely to recent volatility improves risk-adjusted returns
across a range of factors, and \citet{hood25} examine how much of that improvement is
attributable to trend following rather than to volatility management itself.

Table~\ref{tab:paradigms} positions the present study against these paradigms.
This study sits between the two responses. It retains an expected-return term, but
scales it explicitly so that its magnitude is commensurate with the risk term, and it
separates the composition decision from the total-risk decision so that the two can be
evaluated independently.

\begin{table}[!ht]
\centering
\caption{Portfolio construction paradigms and the position of this study}
\label{tab:paradigms}
\small
\begin{tabular}{p{2.4cm}p{2.4cm}p{2.6cm}p{2.6cm}}
\hline
\textbf{Paradigm} & \textbf{Representative work} & \textbf{Basis of allocation} & \textbf{Principal limitation} \\
\hline
Mean--variance optimisation & \citet{markowitz52} & Expected returns and covariance & Extreme sensitivity to estimated inputs \\
\hline
Naive diversification & \citet{demiguel09} & None; equal weights & Ignores risk structure entirely \\
\hline
Shrinkage-based optimisation & \citet{ledoit04, lu25} & Regularised covariance & Still requires expected returns \\
\hline
Risk-based construction & \citet{prado16, asimit25} & Risk model only & Forgoes signal information by design \\
\hline
Volatility management & \citet{moreira17, hood25} & Total risk scaled by recent volatility & Does not address composition \\
\hline
Regime-switching allocation & \citet{angbekaert02, kritzman12, luo26} & State-dependent parameters & Component contributions rarely isolated \\
\hline
Learning-based allocation & \citet{jiang25, ashrafzadeh25, agal25} & Learned policy or predicted returns & Interpretability; frequently single-asset \\
\hline
\textbf{This study} & --- & Regime-conditional signal, plus a relative regime risk budget & Signal layer tested and found not to contribute \\
\hline
\end{tabular}
\end{table}

\subsection{Regime Dependence and Regime-Aware Allocation}

The premise that financial series alternate between statistically distinct states dates
to \citet{hamilton89}, whose Markov-switching formulation remains the reference point.
\citet{angbekaert02} applied regime switching to international allocation and found that
correlations rise in bear states, weakening diversification precisely when it is most
needed---an observation that motivates state-dependent covariance and state-dependent
risk budgets alike. \citet{kritzman12} argued that regime awareness should change the
strategy rather than only the commentary, and \citet{barroso15} demonstrated a
closely related mechanism at the factor level, showing that managing the volatility of
a momentum strategy removes most of its crash risk.

Contemporary work continues along two lines. One identifies regimes with increasingly
flexible statistical machinery: \citet{li25} combine a jump model with allocation
rules, and \citet{luo26} construct dual-regime signals with regime-dependent
parameters. The other embeds state-dependence inside a learning system, where the
regime is implicit in the learned policy \citep{jiang25}. Across both lines, the common
gap is evidential rather than methodological. Regime-aware systems are usually
evaluated against regime-unaware benchmarks in aggregate, so the reader learns that the
system as a whole performs well but not which of its state-dependent components is
responsible. The ablation design in Section~4 is intended to close that gap for the
framework proposed here.

\subsection{Machine Learning and Multimodal Signals}

Machine learning entered this field early and has since become routine. Agent-based
and reinforcement formulations were applied to daily equity trading by \citet{lee07},
metaheuristic-optimised regression to price forecasting by \citet{chou18}, and pattern
mining with fuzzy inference to trading-point detection by \citet{huang20}; each
improved predictive accuracy without addressing what the resulting forecasts imply for a
portfolio. \citet{gu20} provide the benchmark comparative study in asset pricing,
finding consistent gains from flexible models while noting that measured predictability
is small. Applications to allocation
are more recent: \citet{ashrafzadeh25} combine deep learning with stock clustering for
portfolio construction, \citet{agal25} apply learned risk estimates to asset
allocation, \citet{ha26} compare machine-learning portfolios against a fixed
all-weather allocation, and \citet{singsiri26} build portfolios directly on money-flow
indicators of the kind used in this study. Reinforcement learning has been extended from single assets to
multi-period, high-dimensional allocation \citep{jiang25, saijai26}, and
\citet{wang26} couple prediction with an explicit trading strategy.

Fusion of heterogeneous signals is a parallel development. \citet{zhang22} review
decision fusion for equity prediction across a large body of primary studies and find
that fused architectures usually outperform unimodal ones, while also documenting wide
variation in fusion protocol and evaluation. \citet{farimani24} propose an adaptive
multimodal model that weights sources dynamically, and \citet{zong25} use gated
cross-attention to stabilise multimodal fusion for movement prediction. What these
studies establish is that fusion can help prediction. Whether a particular fusion rule
helps a particular allocation decision is a separate question, and one that requires a
component-level test rather than an aggregate comparison.

\subsection{Textual Sentiment and Market-Derived Proxies}

Textual sentiment entered finance through dictionary methods, and
\citet{loughran11} demonstrated that general-purpose sentiment lexicons misclassify
financial language badly enough to require domain-specific alternatives. Transformer
models and large language models have since displaced lexicons: \citet{officioso26}
combine language-model sentiment with market data, \citet{muhammad26} benchmark large
language models for target-based financial sentiment, \citet{mun25} derive investment
signals from language-model sentiment, and \citet{chuang25} compare language models on
financial news for sentiment and price prediction. Collectively these studies indicate
that genuine textual sentiment carries information beyond price history, and
\citet{mazzotta25} show more broadly that alternative data can carry information that
conventional variables do not.

They also indicate what such a study requires: a time-stamped corpus with coverage of
the traded instruments. That requirement is binding. For broad exchange-traded funds
there is no instrument-level news corpus of comparable coverage, which is why studies
of this kind are almost always conducted on individual equities. The present study
therefore uses market-derived proxies for attention and stress rather than textual
sentiment, states this explicitly in Section~3.3, and---rather than assuming the proxies
are adequate---measures their predictive content directly and reports that they are
not.

\subsection{Evaluation Practice}

A final strand concerns whether reported results should be believed.
\citet{harvey15} show that conventional significance thresholds are inadequate once the
number of strategies examined is taken into account, and \citet{valls26} document the
same problem specifically for volatility-based strategies, while \citet{sheppert26}
proposes an objective function designed to resist it. \citet{zhangz25} make the related
point that the model selected as optimal in sample is frequently not the one that
performs best out of it. \citet{ledoit08} provide
robust inference for Sharpe ratio comparisons under non-normal, autocorrelated returns,
and \citet{politis94} supply the stationary block bootstrap used here to construct
interval estimates that respect serial dependence. \citet{newey87} remains the standard
correction for autocorrelated and heteroskedastic test statistics.
\citet{raimundo26} survey the broader intersection of forecasting and portfolio theory
and reach a similar conclusion about the weight placed on point estimates.

This study adopts these practices as design constraints rather than as afterthoughts:
every specification examined is reported, inference is corrected for serial dependence,
interval estimates accompany point estimates, and the comparison against a costless
daily-rebalanced equal-weighted benchmark is deliberately unfavourable to the proposed
framework.
